% Bản slide tự học chi tiết về quá trình cải tiến HRM-PET.
% Biên dịch bằng XeLaTeX, chạy hai lần:
% xelatex -interaction=nonstopmode BAO_CAO_HRM_PET_TU_HOC_CHI_TIET.tex
\documentclass[aspectratio=169,11pt]{beamer}

\usepackage{fontspec}
\usepackage{amsmath,amssymb}
\usepackage{tikz}
\usetikzlibrary{arrows.meta,positioning,calc,fit}
\usepackage{pgfplots}
\pgfplotsset{compat=1.18}
\usepackage{booktabs,tabularx,array,colortbl}
\usepackage{ragged2e}

\setsansfont{TeX Gyre Heros}
\renewcommand{\familydefault}{\sfdefault}

\definecolor{HRMBlue}{HTML}{2563EB}
\definecolor{HRMCyan}{HTML}{DDF4FF}
\definecolor{HRMDark}{HTML}{111827}
\definecolor{HRMGray}{HTML}{6B7280}
\definecolor{HRMLight}{HTML}{F3F4F6}
\definecolor{HRMOrange}{HTML}{F59E0B}
\definecolor{HRMGreen}{HTML}{059669}
\definecolor{HRMRed}{HTML}{DC2626}

\setbeamercolor{normal text}{fg=HRMDark,bg=white}
\setbeamercolor{frametitle}{fg=HRMDark,bg=white}
\setbeamercolor{structure}{fg=HRMBlue}
\setbeamercolor{alerted text}{fg=HRMBlue}
\setbeamercolor{block title}{fg=HRMDark,bg=HRMCyan}
\setbeamercolor{block body}{fg=HRMDark,bg=HRMLight}
\setbeamerfont{frametitle}{size=\LARGE,series=\mdseries}
\setbeamerfont{title}{size=\Huge,series=\mdseries}
\setbeamerfont{subtitle}{size=\large}
\setbeamertemplate{navigation symbols}{}
\setbeamertemplate{itemize item}{\color{HRMBlue}\raise0.15ex\hbox{\small$\bullet$}}
\setbeamertemplate{itemize subitem}{\color{HRMGray}\raise0.15ex\hbox{\scriptsize$\bullet$}}
\setbeamertemplate{frametitle}{%
  \vspace{0.25cm}\hspace{0.1cm}\insertframetitle\par
  \vspace{0.05cm}\hspace{0.1cm}\color{HRMBlue}\rule{1.15cm}{1.5pt}
}
\setbeamertemplate{footline}{%
  \leavevmode
  \hbox{%
    \begin{beamercolorbox}[wd=.86\paperwidth,ht=2.8ex,dp=1.2ex,leftskip=.4cm]{normal text}
      \color{HRMGray}\scriptsize Tài liệu tự học · Cải tiến HRM-PET
    \end{beamercolorbox}%
    \begin{beamercolorbox}[wd=.14\paperwidth,ht=2.8ex,dp=1.2ex,rightskip=.4cm plus1fil]{normal text}
      \color{HRMGray}\scriptsize\insertframenumber/\inserttotalframenumber
    \end{beamercolorbox}%
  }
}

\tikzset{
  flow/.style={draw=HRMGray!35,fill=HRMLight,rounded corners=2pt,
    minimum width=1.9cm,minimum height=0.95cm,align=center,font=\small,inner sep=4pt},
  flownew/.style={flow,draw=HRMBlue,fill=HRMCyan,line width=0.9pt,font=\small\bfseries},
  flowwarn/.style={flow,draw=HRMOrange,fill=HRMOrange!12,line width=0.9pt},
  arrow/.style={-{Latex[length=2.1mm]},draw=HRMGray,line width=0.8pt},
  dashedarrow/.style={arrow,dashed},
  smalllabel/.style={font=\scriptsize\bfseries,text=HRMGray},
  newlabel/.style={font=\scriptsize\bfseries,text=HRMBlue}
}

\newcommand{\sourceNote}[1]{\note{[Sources]\par #1\par[/Sources]}}
\newcommand{\remember}[1]{%
  \vfill
  \begin{beamercolorbox}[wd=\textwidth,sep=0.22cm,rounded=true]{block body}
    \textbf{Ý cần nhớ:} #1
  \end{beamercolorbox}
}
\newcommand{\term}[2]{\textbf{#1}: #2}
\newcommand{\sectionframe}[2]{%
  \begin{frame}[plain]
    \vfill
    {\large\color{HRMBlue}#1}\par\vspace{0.35cm}
    {\Huge\color{HRMDark}#2}\par
    \vfill
  \end{frame}
}

\title{Cải tiến HRM-PET\\Từ cơ chế gốc đến cấu hình Acc@1 88.60\%}
\subtitle{Tài liệu tự học chi tiết · Split-CIFAR100 · 10 task · seed 42}
\author{}
\date{}

\begin{document}

% 1
\begin{frame}[plain]
  \vspace{0.35cm}
  {\large\color{HRMGray}BÁO CÁO TỰ HỌC CHI TIẾT}\par
  \vspace{0.85cm}
  {\usebeamerfont{title}\color{HRMDark}\inserttitle}\par
  \vspace{0.7cm}
  {\usebeamerfont{subtitle}\color{HRMGray}\insertsubtitle}
  \sourceNote{Báo cáo tổng kết, nhật ký thí nghiệm và hai paper nguồn.}
\end{frame}

% 2
\begin{frame}{Đọc tài liệu này theo bốn câu hỏi cố định}
  \vspace{0.35cm}
  \begin{columns}[T,onlytextwidth]
    \begin{column}{0.23\textwidth}
      {\Large\bfseries\color{HRMBlue}1}\par\vspace{0.15cm}
      \textbf{Bản gốc làm gì?}\par
      Hiểu đúng HRM-PET, CTIRD và CRCT trước khi nói đến cải tiến.
    \end{column}
    \begin{column}{0.23\textwidth}
      {\Large\bfseries\color{HRMBlue}2}\par\vspace{0.15cm}
      \textbf{Vấn đề ở đâu?}\par
      Xác định hạn chế từ paper, code thực tế và kết quả chạy.
    \end{column}
    \begin{column}{0.23\textwidth}
      {\Large\bfseries\color{HRMBlue}3}\par\vspace{0.15cm}
      \textbf{Ta sửa thế nào?}\par
      Giải thích CFS, full/boundary replay và task-energy CTIRD.
    \end{column}
    \begin{column}{0.23\textwidth}
      {\Large\bfseries\color{HRMBlue}4}\par\vspace{0.15cm}
      \textbf{Tác dụng ra sao?}\par
      Phân biệt kỳ vọng lý thuyết với kết quả ablation thực tế.
    \end{column}
  \end{columns}
  \remember{Không xem đây là một mô hình hoàn toàn mới. Đây là các cải tiến có chủ đích trên hai vị trí cụ thể của HRM-PET.}
\end{frame}

\sectionframe{PHẦN I}{Hiểu HRM-PET gốc trước khi cải tiến}

% 4
\begin{frame}{Bài toán: học lớp mới mà không quên lớp cũ}
  \begin{columns}[T,onlytextwidth]
    \begin{column}{0.48\textwidth}
      {\large\bfseries Thiết lập Split-CIFAR100}\par\vspace{0.25cm}
      \begin{itemize}\setlength\itemsep{0.35cm}
        \item 100 lớp được chia thành 10 task liên tiếp.
        \item Mỗi task đưa vào 10 lớp mới.
        \item Khi học task sau, dữ liệu ảnh của task trước không còn được dùng trực tiếp.
      \end{itemize}
    \end{column}
    \begin{column}{0.48\textwidth}
      {\large\bfseries Hai khó khăn chính}\par\vspace{0.25cm}
      \begin{itemize}\setlength\itemsep{0.35cm}
        \item \textbf{Forgetting}: cập nhật theo lớp mới làm hiệu năng lớp cũ giảm.
        \item \textbf{Task mismatch}: model có nhiều LoRA; chọn sai LoRA làm dự đoán sai dù backbone tốt.
      \end{itemize}
    \end{column}
  \end{columns}
  \remember{HRM-PET xử lý đồng thời hai chuyện: giữ tri thức qua các task và sửa việc chọn nhầm tham số LoRA.}
  \sourceNote{HRM-PET paper, problem formulation and experiments.}
\end{frame}

% 5
\begin{frame}{Từ điển nhanh: các thành phần cần nhớ}
  \small
  \begin{tabularx}{\textwidth}{>{\bfseries}p{2.5cm} X}
    \toprule
    Thuật ngữ & Ý nghĩa trong hệ thống \\
    \midrule
    ViT backbone & Mô hình thị giác tiền huấn luyện, phần lớn được đóng băng. \\
    LoRA/PET & Khối tham số nhỏ cho từng task; giúp thích nghi mà không fine-tune toàn bộ ViT. \\
    TII & Task Identity Inference: dự đoán mẫu thuộc task nào để chọn LoRA tương ứng. \\
    DRM/CRM & Hai cơ chế re-matching ở lúc suy luận để sửa trường hợp chọn nhầm task/LoRA. \\
    CTIRD & Distillation quan hệ giữa các mẫu qua LoRA hiện tại và LoRA cũ. \\
    Feature replay & Feature tổng hợp đại diện lớp cũ; không phải ảnh gốc. \\
    CRCT & Classifier correction: học lại classifier bằng feature replay để cân bằng lớp cũ/mới. \\
    \bottomrule
  \end{tabularx}
  \remember{CFS tác động lên feature replay/CRCT; task energy tác động lên CTIRD. DRM/CRM và backbone được giữ nguyên.}
\end{frame}

% 6
\begin{frame}{HRM-PET gốc khi suy luận: chọn LoRA rồi re-match}
  \centering
  \vspace{0.35cm}
  \begin{tikzpicture}[node distance=0.38cm]
    \node[flow] (a) {Ảnh kiểm tra\\$x$};
    \node[flow] (b) [right=of a] {TII dự đoán\\task $\hat t_f$};
    \node[flow] (c) [right=of b] {Gắn LoRA\\$p_{\hat t_f}$};
    \node[flow] (d) [right=of c] {Classifier\\dự đoán lớp};
    \node[flownew] (e) [right=of d] {DRM/CRM\\re-match};
    \node[flow] (f) [right=of e] {Kết quả\\cuối};
    \foreach \x/\y in {a/b,b/c,c/d,d/e,e/f} \draw[arrow] (\x)--(\y);
  \end{tikzpicture}

  \vspace{0.7cm}
  \begin{columns}[T,onlytextwidth]
    \begin{column}{0.48\textwidth}
      \term{DRM}{dùng nhãn lớp dự đoán để suy ra lại task, sau đó thử LoRA tương ứng.}
    \end{column}
    \begin{column}{0.48\textwidth}
      \term{CRM}{khi dự đoán thiếu tin cậy, so sánh confidence của một số LoRA ứng viên.}
    \end{column}
  \end{columns}
  \remember{Các cải tiến của chúng ta không thay luồng suy luận này.}
  \sourceNote{HRM-PET paper, Section 3.2: Direct and Confidence-based Re-matching.}
\end{frame}

% 7
\begin{frame}{HRM-PET gốc khi huấn luyện: hai nhánh bù dữ liệu cũ}
  \centering
  \vspace{0.25cm}
  \begin{tikzpicture}[node distance=0.42cm]
    \node[flow] (a) {Dữ liệu\\task mới};
    \node[flow] (b) [right=of a] {ViT + LoRA\\task hiện tại};
    \node[flownew] (c) [above right=0.25cm and 0.55cm of b] {CTIRD\\giữ quan hệ};
    \node[flownew] (d) [below right=0.25cm and 0.55cm of b] {Lưu thống kê\\feature/lớp};
    \node[flow] (e) [right=1.0cm of d] {Gaussian\\feature replay};
    \node[flow] (f) [right=of e] {CRCT sửa\\classifier};
    \draw[arrow] (a)--(b);
    \draw[arrow] (b)--(c);
    \draw[arrow] (b)--(d);
    \draw[arrow] (d)--(e);
    \draw[arrow] (e)--(f);
  \end{tikzpicture}

  \vspace{0.55cm}
  \begin{columns}[T,onlytextwidth]
    \begin{column}{0.48\textwidth}
      \textbf{Nhánh CTIRD}\par
      Ép LoRA mới học quan hệ giữa các mẫu tương thích với tri thức chia sẻ từ LoRA cũ.
    \end{column}
    \begin{column}{0.48\textwidth}
      \textbf{Nhánh replay + CRCT}\par
      Mô phỏng feature lớp cũ rồi dùng chúng để cân lại classifier sau mỗi task.
    \end{column}
  \end{columns}
  \remember{Hai nhánh độc lập về vai trò: CTIRD định hình feature; CRCT sửa quyết định phân lớp.}
\end{frame}

% 8
\begin{frame}{CTIRD gốc: giữ cấu trúc quan hệ giữa các mẫu}
  \begin{columns}[T,onlytextwidth]
    \begin{column}{0.52\textwidth}
      Với batch $\{x_i\}_{i=1}^{B}$, LoRA hiện tại tạo feature $e_t(x_i)$ và ma trận quan hệ:
      \[
        W_t^{ij}=\langle \bar e_t(x_i),\bar e_t(x_j)\rangle.
      \]
      Một LoRA cũ $p_k$ tạo ma trận $W_k$ tương tự. CTIRD căn chỉnh các phân phối quan hệ:
      \[
        \mathcal L_{\mathrm{CTIRD}}
        =\sum_{k\in\mathcal S}\mathrm{KL}
        \bigl(\mathrm{Norm}(W_t)\,\|\,\mathrm{Norm}(W_k)\bigr).
      \]
    \end{column}
    \begin{column}{0.44\textwidth}
      \begin{block}{Trực giác}
        Không bắt feature mới giống hệt feature cũ. Chỉ yêu cầu các mẫu gần/xa nhau theo cấu trúc tương tự.
      \end{block}
      \vspace{0.25cm}
      \begin{block}{Mục tiêu}
        Học tri thức dùng chung giữa task, để khi gắn nhầm LoRA model vẫn giữ được vùng chú ý và quan hệ hữu ích.
      \end{block}
    \end{column}
  \end{columns}
  \remember{CTIRD mạnh hay nhiễu phụ thuộc trực tiếp vào việc tập source LoRA $\mathcal S$ có liên quan hay không.}
  \sourceNote{HRM-PET paper, Section 3.3 and Eq. 8--11.}
\end{frame}

% 9
\begin{frame}{Feature replay gốc: lưu phân phối, không lưu ảnh}
  Với feature lớp $c$, code lưu trung bình và covariance:
  \[
    \mu_c=\frac{1}{N_c}\sum_i z_i,
    \qquad
    \Sigma_c=\frac{1}{N_c-1}\sum_i(z_i-\mu_c)(z_i-\mu_c)^\top.
  \]
  Khi correction, feature giả được lấy mẫu:
  \[
    \tilde z_c\sim\mathcal N(\mu_c,\Sigma_c).
  \]

  \begin{columns}[T,onlytextwidth]
    \begin{column}{0.48\textwidth}
      \begin{block}{Ưu điểm}
        Nhẹ hơn lưu ảnh; không cần chạy backbone cho replay; phù hợp data-free/replay-free ở cấp ảnh.
      \end{block}
    \end{column}
    \begin{column}{0.48\textwidth}
      \begin{block}{Giới hạn}
        Mean/covariance chỉ mô tả thô phân phối. Nhiều mẫu Gaussian có thể gần nhau và cung cấp gradient lặp lại.
      \end{block}
    \end{column}
  \end{columns}
  \remember{CFS không thay Gaussian; nó chọn tốt hơn trong số các mẫu Gaussian có thể sinh.}
  \sourceNote{HRM-PET code: hide\_tii\_engine.py and hrm\_lora\_wtp\_and\_tap\_engine.py.}
\end{frame}

% 10
\begin{frame}{CRCT gốc: học lại classifier bằng feature tổng hợp}
  \centering
  \begin{tikzpicture}[node distance=0.55cm]
    \node[flow] (a) {Thống kê\\từng lớp};
    \node[flow] (b) [right=of a] {Sinh $5\times$ batch\\mẫu/lớp};
    \node[flow] (c) [right=of b] {Trộn tất cả\\lớp đã thấy};
    \node[flow] (d) [right=of c] {Chỉ mở\\classifier head};
    \node[flow] (e) [right=of d] {Cross-entropy\\correction};
    \foreach \x/\y in {a/b,b/c,c/d,d/e} \draw[arrow] (\x)--(\y);
  \end{tikzpicture}

  \vspace{0.6cm}
  \begin{itemize}\setlength\itemsep{0.3cm}
    \item Mỗi lớp sinh \texttt{num\_sampled\_pcls = batch\_size $\times$ 5} feature.
    \item Backbone/LoRA không cần cập nhật ở bước này; CRCT tập trung sửa classifier bị thiên về lớp mới.
    \item Code cũ đặt số vòng replay theo \texttt{crct\_num}, không bảo đảm đi qua toàn bộ tensor replay đã sinh.
  \end{itemize}
  \remember{Full replay sau này sửa đúng điểm cuối: số vòng phải phủ hết số batch thực tế trong replay tensor.}
\end{frame}

% 11
\begin{frame}{Ba nút thắt dẫn đến toàn bộ phần cải tiến}
  \begin{columns}[T,onlytextwidth]
    \begin{column}{0.31\textwidth}
      {\Large\bfseries\color{HRMBlue}1}\par
      \textbf{Replay chưa chọn lọc}\par\vspace{0.25cm}
      Gaussian sample hợp lệ nhưng có thể trùng nhau, ít đại diện cho chi tiết phân phối.
    \end{column}
    \begin{column}{0.31\textwidth}
      {\Large\bfseries\color{HRMBlue}2}\par
      \textbf{Replay chưa dùng hết}\par\vspace{0.25cm}
      CRCT đã sinh nhiều feature nhưng số vòng cũ có thể bỏ lại phần cuối.
    \end{column}
    \begin{column}{0.31\textwidth}
      {\Large\bfseries\color{HRMBlue}3}\par
      \textbf{CTIRD chọn nguồn nhiễu}\par\vspace{0.25cm}
      Code baseline chọn class logit thấp rồi map sang task; source có thể ít liên quan hoặc trùng.
    \end{column}
  \end{columns}
  \remember{CFS giải quyết nút 1; full/boundary replay giải quyết nút 2; task-energy CTIRD giải quyết nút 3.}
\end{frame}

\sectionframe{PHẦN II}{Cải thiện feature replay và CRCT}

% 13
\begin{frame}{CFS trong paper nguồn xuất phát từ vấn đề gì?}
  \begin{columns}[T,onlytextwidth]
    \begin{column}{0.48\textwidth}
      {\large\bfseries Pipeline paper nguồn}\par\vspace{0.3cm}
      \begin{enumerate}
        \item Mô hình hóa feature thật theo Gaussian.
        \item Dùng CFS chọn feature mục tiêu đa dạng.
        \item Dùng model inversion sinh ảnh có feature khớp mục tiêu.
      \end{enumerate}
    \end{column}
    \begin{column}{0.48\textwidth}
      {\large\bfseries Ý tưởng cốt lõi}\par\vspace{0.3cm}
      Gaussian chỉ giữ mean/variance. Một MLP contrastive được học từ feature thật để tạo không gian mà diversity dễ đo hơn; sau đó chọn candidate trải đều trong không gian đó.
    \end{column}
  \end{columns}
  \remember{Ta chuyển giao phần ``contrastive feature selection'', không bê nguyên model inversion hay sinh ảnh.}
  \sourceNote{Model Inversion with Layer-wise Optimization, Section 3.3.}
\end{frame}

% 14
\begin{frame}{Bước 1 của CFS: học một MLP contrastive cho từng lớp}
  Với feature thật $z_i$ của lớp $c$, MLP hai tầng tạo embedding chuẩn hóa:
  \[
    u_i=\frac{f_{\mathrm{CFS}}(z_i)}{\|f_{\mathrm{CFS}}(z_i)\|_2}.
  \]
  Code tối thiểu hóa mức tương đồng với các feature khác trong batch:
  \[
    \mathcal L_{\mathrm{CFS}}
    =\frac{1}{B}\sum_i
    \log\sum_{j\ne i}\exp\left(\frac{u_i^\top u_j}{\tau}\right).
  \]

  \begin{columns}[T,onlytextwidth]
    \begin{column}{0.48\textwidth}
      \textbf{Khi tối thiểu hóa loss}\par
      Embedding của các feature trong cùng lớp bị đẩy ra xa nhau trên hypersphere.
    \end{column}
    \begin{column}{0.48\textwidth}
      \textbf{Thông số triển khai}\par
      MLP hai tầng, LeakyReLU, hidden tối đa 512, SGD; có thể giới hạn tối đa 1024 feature/lớp.
    \end{column}
  \end{columns}
  \remember{MLP không phải classifier. Nó chỉ tạo một thước đo diversity tốt hơn để chọn replay.}
  \sourceNote{utils.py: CFSContrastiveMLP and train\_cfs\_model; paper CFS Eq. 6.}
\end{frame}

% 15
\begin{frame}{Bước 2 của CFS: chọn tham lam các candidate ít giống nhau}
  \centering
  \begin{tikzpicture}[node distance=0.42cm]
    \node[flow] (a) {Sinh $M$ candidate\\từ Gaussian};
    \node[flow] (b) [right=of a] {Chiếu candidate\\qua MLP CFS};
    \node[flow] (c) [right=of b] {Khởi tạo\\một mẫu};
    \node[flownew] (d) [right=of c] {Chọn mẫu có\\similarity TB nhỏ nhất};
    \node[flownew] (e) [right=of d] {Lặp đến đủ\\số replay};
    \foreach \x/\y in {a/b,b/c,c/d,d/e} \draw[arrow] (\x)--(\y);
  \end{tikzpicture}

  \vspace{0.55cm}
  Nếu tập đã chọn là $S$, candidate tiếp theo thỏa:
  \[
    z^*=\arg\min_{z\in C\setminus S}
    \frac{1}{|S|}\sum_{s\in S}
    \exp\left(\frac{f(z)^\top f(s)}{\tau}\right).
  \]
  \begin{itemize}\setlength\itemsep{0.2cm}
    \item \texttt{candidate\_multiplier} điều khiển số candidate so với số replay cần chọn.
    \item Bản tốt nhất dùng khởi tạo ngẫu nhiên; mean initialization và distribution filter không cải thiện.
  \end{itemize}
  \remember{CFS không tạo thêm nhãn hay lớp mới; nó chỉ chọn subset đa dạng hơn trong cùng phân phối lớp.}
\end{frame}

% 16
\begin{frame}{CFS được gắn vào HRM-PET ở đâu?}
  \centering
  \vspace{0.35cm}
  \begin{tikzpicture}[node distance=0.38cm]
    \node[flow] (a) {Feature thật\\theo lớp};
    \node[flow] (b) [right=of a] {Mean +\\covariance};
    \node[flow] (c) [right=of b] {Nhiều Gaussian\\candidate};
    \node[flownew] (d) [right=of c] {CFS chọn\\subset đa dạng};
    \node[flow] (e) [right=of d] {Replay tensor\\có nhãn lớp};
    \node[flow] (f) [right=of e] {CRCT sửa\\classifier};
    \foreach \x/\y in {a/b,b/c,c/d,d/e,e/f} \draw[arrow] (\x)--(\y);
  \end{tikzpicture}

  \vspace{0.65cm}
  \begin{columns}[T,onlytextwidth]
    \begin{column}{0.48\textwidth}
      \textbf{Áp dụng trong cả hai engine}\par
      TII engine và LoRA engine đều có bước thống kê feature và classifier correction, nên cùng gọi hàm sampling CFS.
    \end{column}
    \begin{column}{0.48\textwidth}
      \textbf{Không thay đổi}\par
      Không đổi ViT, LoRA, dữ liệu ảnh, DRM/CRM và không thêm chi phí CFS khi inference.
    \end{column}
  \end{columns}
  \remember{Đây là ``áp dụng ý tưởng CFS vào feature replay'', không nên trình bày là sao chép nguyên pipeline paper nguồn.}
\end{frame}

% 17
\begin{frame}{Full replay: sửa việc CRCT sinh nhiều nhưng dùng chưa hết}
  \begin{columns}[T,onlytextwidth]
    \begin{column}{0.47\textwidth}
      {\large\bfseries Code cũ}\par\vspace{0.25cm}
      \[
        \texttt{replay\_iterations}=\texttt{crct\_num}.
      \]
      Số vòng gắn với số lớp cũ, trong khi tensor replay chứa feature của toàn bộ lớp đã thấy. Phần cuối tensor có thể không được đưa qua optimizer.
    \end{column}
    \begin{column}{0.47\textwidth}
      {\large\bfseries Sau cải tiến}\par\vspace{0.25cm}
      \[
        \texttt{replay\_iterations}
        =\left\lceil\frac{N_{\mathrm{replay}}}{N_{\mathrm{batch}}}\right\rceil.
      \]
      Flag \texttt{--crct\_use\_all\_samples} bảo đảm mỗi epoch correction đi qua toàn bộ replay tensor.
    \end{column}
  \end{columns}
  \vspace{0.4cm}
  \begin{block}{Tác dụng thực tế}
    Đây là bước tăng rõ đầu tiên: Acc@1 từ 87.88 lên 88.13 và loss giảm từ 0.5222 xuống 0.4696.
  \end{block}
  \remember{CFS tăng chất lượng replay; full replay bảo đảm chất lượng đó thực sự được optimizer sử dụng.}
\end{frame}

% 18
\begin{frame}{Boundary-aware replay: thêm một tỷ lệ nhỏ mẫu khó}
  \small
  Với candidate $z$ của lớp mục tiêu $c$, classifier cho margin:
  \[
    m(z)=\ell_c(z)-\max_{j\ne c}\ell_j(z).
  \]
  Mẫu gần biên có $|m(z)|$ nhỏ. Trước khi chọn, code bỏ phần đuôi Gaussian quá xa bằng điểm Mahalanobis xấp xỉ:
  \[
    d(z)=\operatorname{mean}\left(\frac{(z-\mu_c)^2}{\operatorname{diag}(\Sigma_c)}\right).
  \]

  \begin{enumerate}\setlength\itemsep{0.16cm}
    \item Giữ $75\%$ replay theo CFS diversity.
    \item Sinh thêm candidate từ Gaussian.
    \item Giữ candidate còn nằm trong vùng phân phối lớp.
    \item Chọn $25\%$ có margin gần 0 nhất để classifier học vùng dễ nhầm.
  \end{enumerate}
  \remember{Boundary replay không thay CFS; nó bổ sung độ khó có kiểm soát. Tỷ lệ lớn quá có thể làm mất diversity.}
  \sourceNote{utils.py: sample\_boundary\_aware\_cfs\_features; báo cáo ablation boundary ratio.}
\end{frame}

% 19
\begin{frame}{Ba thay đổi replay bổ sung nhau, không thay thế nhau}
  \begin{columns}[T,onlytextwidth]
    \begin{column}{0.31\textwidth}
      {\Large\bfseries CFS}\par\vspace{0.2cm}
      \textbf{Câu hỏi:} chọn mẫu nào?\par\vspace{0.2cm}
      Ưu tiên feature ít trùng và đại diện nhiều vùng của phân phối.
    \end{column}
    \begin{column}{0.31\textwidth}
      {\Large\bfseries Full replay}\par\vspace{0.2cm}
      \textbf{Câu hỏi:} dùng bao nhiêu?\par\vspace{0.2cm}
      Đi qua toàn bộ feature đã sinh trong mỗi epoch CRCT.
    \end{column}
    \begin{column}{0.31\textwidth}
      {\Large\bfseries Boundary}\par\vspace{0.2cm}
      \textbf{Câu hỏi:} có đủ mẫu khó chưa?\par\vspace{0.2cm}
      Dành một phần replay cho vùng classifier dễ nhầm.
    \end{column}
  \end{columns}
  \remember{Chất lượng, mức sử dụng và độ khó là ba trục khác nhau; vì vậy ba cải tiến có thể cùng bật.}
\end{frame}

\sectionframe{PHẦN III}{Sửa cách CTIRD chọn và dùng source task}

% 21
\begin{frame}{Ý định trong paper đúng, nhưng code baseline chọn khác}
  \begin{columns}[T,onlytextwidth]
    \begin{column}{0.48\textwidth}
      {\large\bfseries Paper mô tả}\par\vspace{0.25cm}
      Dùng task classifier để lấy top-K task identity có confidence cao cho mỗi mẫu; chạy các LoRA đó để tạo relation target cho CTIRD.
    \end{column}
    \begin{column}{0.48\textwidth}
      {\large\bfseries Code baseline thực tế}\par\vspace{0.25cm}
      \texttt{torch.topk(..., largest=False)} trên class logits cũ, rồi map từng class sang task. Hệ quả:
      \begin{itemize}
        \item chọn phía logit thấp;
        \item nhiều class có thể map trùng một task;
        \item source LoRA có thể ít liên quan.
      \end{itemize}
    \end{column}
  \end{columns}
  \remember{Ta không bỏ CTIRD. Ta sửa nguồn distillation để CTIRD thực hiện đúng mục tiêu hơn.}
  \sourceNote{HRM-PET paper, Section 3.3; engine code legacy branch in get\_old\_features/train\_one\_epoch.}
\end{frame}

% 22
\begin{frame}{Task energy: xếp hạng trực tiếp ở cấp task}
  Với task $t$ gồm tập class $\mathcal C_t$, energy được gộp bằng log-sum-exp:
  \[
    E_t(x)=T\log\sum_{c\in\mathcal C_t}
    \exp\left(\frac{\ell_c(x)}{T}\right).
  \]

  \begin{columns}[T,onlytextwidth]
    \begin{column}{0.48\textwidth}
      \textbf{Vì sao log-sum-exp?}\par
      Nó tổng hợp bằng chứng của tất cả class trong task nhưng vẫn nhấn mạnh class có logit lớn. Nhiệt độ $T$ điều khiển độ sắc.
    \end{column}
    \begin{column}{0.48\textwidth}
      \textbf{Cách chọn}\par
      Tính một score cho mỗi task cũ, sau đó lấy top-K với \texttt{largest=True}. Mỗi phần tử top-K là một task riêng biệt.
    \end{column}
  \end{columns}
  \vspace{0.35cm}
  Cấu hình tốt nhất: $K=5$, task temperature $T=0.1$.
  \remember{Thay vì hỏi ``class cũ nào có logit thấp?'', code mới hỏi ``task cũ nào giải thích mẫu này tốt nhất?''.}
\end{frame}

% 23
\begin{frame}{Energy weighting: task liên quan hơn nhận lực distillation lớn hơn}
  Nếu $E_{ik}$ là energy của source task hạng $k$ cho mẫu $i$:
  \[
    q_{ik}=\operatorname{softmax}_k(E_{ik}/\tau_w).
  \]
  Để tránh một task chiếm toàn bộ trọng số, trộn thêm uniform floor $\rho$:
  \[
    \tilde q_{ik}=(1-\rho)q_{ik}+\rho/K.
  \]
  Code lấy trung bình theo batch và nhân $K$:
  \[
    w_k=K\cdot\frac{1}{B}\sum_i\tilde q_{ik},
    \qquad \sum_k w_k\approx K.
  \]
  Loss cuối:
  \[
    \mathcal L=\mathcal L_{CE}+\lambda_{CT}\sum_{k=1}^{K}w_k\mathcal L_{CTIRD}^{(k)}.
  \]
  \remember{Tổng lực CTIRD gần như giữ nguyên; ta chỉ chuyển lực từ source ít liên quan sang source liên quan hơn.}
\end{frame}

% 24
\begin{frame}{CFS và task-energy CTIRD không xung đột nhau}
  \centering
  \begin{tikzpicture}
    \node[flownew,minimum width=4.8cm,minimum height=1.25cm] (a) {CFS + full/boundary replay\\Sửa dữ liệu dùng để correction classifier};
    \node[flownew,minimum width=4.8cm,minimum height=1.25cm,right=1.2cm of a] (b) {Task-energy CTIRD\\Sửa nguồn quan hệ dùng khi train LoRA};
    \node[flow,minimum width=5.2cm,minimum height=1.2cm,below=1.0cm of $(a)!0.5!(b)$] (c) {Feature tốt hơn + classifier cân bằng hơn\\$\Rightarrow$ dự đoán cuối tốt hơn};
    \draw[arrow] (a)--(c);
    \draw[arrow] (b)--(c);
  \end{tikzpicture}

  \vspace{0.45cm}
  \begin{itemize}\setlength\itemsep{0.25cm}
    \item CFS chạy ở pha thống kê/replay; CTIRD chạy khi học LoRA trên ảnh task hiện tại.
    \item Hai nhánh có thể tác động gián tiếp qua chất lượng feature, nhưng không tối ưu cùng một tensor hay cùng một loss.
    \item Xung đột chỉ xuất hiện nếu replay/semantic ép feature lệch quá mạnh, không phải do CFS và CTIRD vốn đối nghịch.
  \end{itemize}
  \remember{Một nhánh sửa ``dữ liệu nhớ'', nhánh kia sửa ``nguồn tri thức''.}
\end{frame}

\sectionframe{PHẦN IV}{Semantic-aware: đã thử nhưng không giữ trong bản cuối}

% 26
\begin{frame}{Semantic-aware projection trong paper nguồn hoạt động thế nào?}
  Paper dùng CLIP, nơi text và image feature cùng nằm trên hypersphere. Với lớp cũ $c$ và lớp mới $d$:
  \[
    F_t(t_d)=R_{c,d}F_t(t_c),
    \qquad
    \tilde z_d=R_{c,d}z_c.
  \]
  Trong đó $R_{c,d}$ là phép quay đưa text embedding lớp $c$ sang lớp $d$, rồi áp dụng cùng phép quay lên image feature.

  \begin{block}{Tại sao hấp dẫn?}
    Nếu chỉ có feature lớp cũ, semantic class name có thể gợi ý cách tạo pseudo-feature mang nghĩa của lớp mới hoặc cách xác định lớp liên quan.
  \end{block}
  \remember{Giả định quan trọng của paper là text và image feature đã được CLIP căn chỉnh. HRM-PET không tự động có giả định này.}
  \sourceNote{Model Inversion with Layer-wise Optimization, Section 3.4 and Appendix I.3.}
\end{frame}

% 27
\begin{frame}{Ta đã thử semantic ở hai vị trí}
  \begin{columns}[T,onlytextwidth]
    \begin{column}{0.48\textwidth}
      {\large\bfseries 1 · Semantic cho replay/CRCT}\par\vspace{0.25cm}
      \begin{itemize}
        \item mean shift và paper-style rotation;
        \item covariance transfer;
        \item top-5 lớp semantic gần nhất;
        \item filter theo phân phối target;
        \item feature adapter căn semantic vào HRM mean.
      \end{itemize}
    \end{column}
    \begin{column}{0.48\textwidth}
      {\large\bfseries 2 · Semantic cho CTIRD}\par\vspace{0.25cm}
      \begin{itemize}
        \item semantic top-k;
        \item adaptive semantic gate;
        \item điều chỉnh relation target theo quan hệ nhãn;
        \item thử backend hash và CLIP.
      \end{itemize}
    \end{column}
  \end{columns}
  \remember{Semantic đã được triển khai thành một nhánh ablation đầy đủ, không phải chỉ là đề xuất trên giấy.}
\end{frame}

% 28
\begin{frame}{Vì sao semantic chưa được đưa vào cấu hình tốt nhất?}
  \begin{columns}[T,onlytextwidth]
    \begin{column}{0.46\textwidth}
      {\large\bfseries Quan sát thực nghiệm}\par\vspace{0.25cm}
      \begin{itemize}\setlength\itemsep{0.25cm}
        \item Một số bản giảm loss hoặc forgetting nhẹ.
        \item Acc@1 dao động quanh 87.14--87.86 tùy biến thể.
        \item Không vượt CFS-only một cách ổn định.
      \end{itemize}
    \end{column}
    \begin{column}{0.50\textwidth}
      {\large\bfseries Giải thích hợp lý}\par\vspace{0.25cm}
      \begin{itemize}\setlength\itemsep{0.25cm}
        \item HRM ViT feature và text embedding không cùng không gian như CLIP.
        \item Class name không đủ mô tả biến thiên thị giác trong CIFAR100.
        \item Prior semantic mạnh có thể kéo replay/relation target lệch khỏi phân phối thật.
      \end{itemize}
    \end{column}
  \end{columns}
  \remember{Quyết định tắt semantic là kết luận từ ablation, không có nghĩa semantic vô ích trong mọi thiết lập.}
\end{frame}

\sectionframe{PHẦN V}{Thiết kế thí nghiệm và kết quả}

% 30
\begin{frame}{Thiết lập đánh giá và nguyên tắc so sánh}
  \begin{columns}[T,onlytextwidth]
    \begin{column}{0.48\textwidth}
      {\large\bfseries Thiết lập chính}\par\vspace{0.2cm}
      \begin{itemize}\setlength\itemsep{0.22cm}
        \item Dataset: Split-CIFAR100, 10 task.
        \item Backbone: ViT-B/16, LoRA rank 5.
        \item Seed: 42; máy chính RTX 4090 24GB.
        \item Cùng TII checkpoint cho chuỗi ablation LoRA/CRCT.
      \end{itemize}
    \end{column}
    \begin{column}{0.48\textwidth}
      {\large\bfseries Mức độ tin cậy}\par\vspace{0.2cm}
      \begin{itemize}\setlength\itemsep{0.22cm}
        \item Baseline gốc đã được chạy lại trên cùng RTX 4090: Acc@1 87.81.
        \item Log Kaggle 86.90 chỉ còn dùng như kết quả lịch sử tham chiếu.
        \item Mới có một seed nên chưa được tuyên bố ý nghĩa thống kê.
      \end{itemize}
    \end{column}
  \end{columns}
  \remember{So sánh chính hiện nay là bản gốc 87.81 với bản cải tiến 88.60 trên cùng máy: tăng 0.79 điểm Acc@1.}
\end{frame}

% 31
\begin{frame}{Các chỉ số trong log nên được đọc thế nào?}
  \small
  \begin{tabularx}{\textwidth}{>{\bfseries}p{2.25cm} X p{4.2cm}}
    \toprule
    Chỉ số & Ý nghĩa & Hướng mong muốn \\
    \midrule
    Acc@1 & Tỷ lệ nhãn đúng đứng đầu dự đoán lớp cuối. & Càng cao càng tốt; đây là mục tiêu chính. \\
    Acc@5 & Nhãn đúng nằm trong năm dự đoán đầu. & Càng cao càng tốt; ít nhạy hơn Acc@1. \\
    Acc@task & Chỉ số task-aware/task matching do pipeline HRM-PET báo cáo. & Càng cao càng tốt, phản ánh chất lượng xác định/re-match task. \\
    Loss & Cross-entropy đánh giá cuối. & Càng thấp càng tốt nếu accuracy không giảm. \\
    Forgetting & Mức giảm trung bình của task cũ so với thời điểm tốt nhất của chúng. & Càng thấp càng tốt. \\
    Backward & Ảnh hưởng của việc học task sau lên task trước. Giá trị âm biểu thị suy giảm. & Gần 0 hoặc dương tốt hơn. \\
    \bottomrule
  \end{tabularx}
  \remember{Không nên chỉ nhìn Acc@1; cấu hình cuối tăng accuracy nhưng forgetting vẫn là trade-off cần xử lý.}
\end{frame}

\begin{frame}{Baseline gốc chạy lại: đây là mốc so sánh chính}
  \begin{columns}[T,onlytextwidth]
    \begin{column}{0.47\textwidth}
      {\large\bfseries HRM-PET gốc, task 10}\par\vspace{0.2cm}
      \begin{tabular}{lr}
        Acc@task & 88.01 \\
        Acc@1 & 87.81 \\
        Acc@5 & 98.05 \\
        Loss & 0.5228 \\
        Forgetting & 3.5889 \\
        Backward & $-3.2556$ \\
      \end{tabular}
    \end{column}
    \begin{column}{0.49\textwidth}
      {\large\bfseries Cách hiểu}\par\vspace{0.2cm}
      \begin{itemize}\setlength\itemsep{0.24cm}
        \item Đây là dòng đánh giá cuối sau classifier correction, không phải kết quả trung gian 78.27 trước đó.
        \item Vì được chạy lại trên cùng máy, dataset và seed, mốc này phù hợp hơn log Kaggle để đánh giá cải tiến.
        \item Baseline đã khá mạnh; mức tăng tiếp theo vì vậy cần được đọc theo cả accuracy, loss và forgetting.
      \end{itemize}
    \end{column}
  \end{columns}
  \remember{Từ đây, mọi tuyên bố \emph{tăng bao nhiêu} trong báo cáo dùng 87.81 làm Acc@1 gốc.}
\end{frame}

% 32
\begin{frame}{Ablation: mỗi thay đổi đóng góp từng phần}
  \centering
  \renewcommand{\arraystretch}{1.23}
  \setlength{\tabcolsep}{4pt}
  \small
  \begin{tabularx}{\textwidth}{>{\RaggedRight\arraybackslash}p{2.45cm} X >{\centering\arraybackslash}p{1.2cm} >{\centering\arraybackslash}p{1.2cm} >{\centering\arraybackslash}p{1.25cm}}
    \toprule
    \textbf{Giai đoạn} & \textbf{Thay đổi chính} & \textbf{Acc@task} & \textbf{Acc@1} & \textbf{Loss} \\
    \midrule
    HRM-PET gốc (Kaggle) & Gaussian replay + CTIRD baseline & 87.05 & 86.90 & 0.5860 \\
    HRM-PET gốc (4090) & Chạy lại code gốc cùng máy & 88.01 & 87.81 & 0.5228 \\
    CFS-only (4090) & Chọn replay đa dạng hơn & 87.62 & 87.88 & 0.5222 \\
    + Full CRCT & Dùng toàn bộ replay & 88.01 & 88.13 & 0.4696 \\
    + Boundary + task energy & Mẫu khó + top-5 task liên quan & 88.05 & 88.31 & 0.4675 \\
    + Energy weighting & Gán trọng số theo source energy & 88.30 & 88.56 & 0.4627 \\
    \rowcolor{HRMCyan} + CRCT tuning & \texttt{ca\_lr=0.0045} & \textbf{88.32} & \textbf{88.60} & \textbf{0.4607} \\
    \bottomrule
  \end{tabularx}
  \vspace{0.2cm}
  {\scriptsize\color{HRMGray}Hàng Kaggle chỉ tham chiếu lịch sử; hàng HRM-PET gốc (4090) là baseline chính cho so sánh cuối.}
\end{frame}

% 33
\begin{frame}{Cấu hình tốt nhất hiện tại gồm những gì?}
  \small
  \begin{columns}[T,onlytextwidth]
    \begin{column}{0.48\textwidth}
      {\large\bfseries LoRA và CTIRD}\par\vspace{0.18cm}
      \begin{tabular}{ll}
        LoRA epochs & 10 \\
        LoRA lr / rank & 0.03 / 5 \\
        Momentum & 0.4 \\
        CTIRD $\lambda$ (con) & 0.2 \\
        K source tasks & 5 \\
        Selection & task energy \\
        Task temperature & 0.1 \\
        Weighting / floor & energy / 0.2 \\
      \end{tabular}
    \end{column}
    \begin{column}{0.48\textwidth}
      {\large\bfseries Replay và CRCT}\par\vspace{0.18cm}
      \begin{tabular}{ll}
        CFS epochs & 20 \\
        Candidate multiplier & 2 \\
        Full replay & bật \\
        Boundary ratio & 0.25 \\
        CRCT epochs & 3 \\
        CRCT lr & 0.0045 \\
        Semantic & tắt \\
        Balanced batches & tắt \\
      \end{tabular}
    \end{column}
  \end{columns}
  \remember{Cấu hình cuối ưu tiên Acc@1. Nếu ưu tiên cân bằng forgetting/loss, \texttt{ca\_lr=0.004} là phương án gần tương đương.}
\end{frame}

% 34
\begin{frame}{Kết quả cuối: accuracy tăng, loss giảm}
  \begin{columns}[T,onlytextwidth]
    \begin{column}{0.58\textwidth}
      \centering
      \begin{tikzpicture}
        \begin{axis}[
          ybar,ymin=84.5,ymax=100,width=8.2cm,height=5.2cm,bar width=12pt,
          symbolic x coords={Acc@task,Acc@1,Acc@5},xtick=data,
          ylabel={Điểm (\%)},tick label style={font=\scriptsize},label style={font=\scriptsize},
          nodes near coords,nodes near coords style={font=\tiny},
          ymajorgrids=true,grid style={draw=HRMGray!18},
          legend style={at={(0.5,-0.18)},anchor=north,legend columns=2,draw=none,font=\scriptsize}]
          \addplot[fill=HRMGray!45,draw=none] coordinates {(Acc@task,88.01) (Acc@1,87.81) (Acc@5,98.05)};
          \addplot[fill=HRMBlue!85,draw=none] coordinates {(Acc@task,88.32) (Acc@1,88.60) (Acc@5,98.07)};
          \legend{HRM-PET gốc cùng máy,Cấu hình cuối}
        \end{axis}
      \end{tikzpicture}
    \end{column}
    \begin{column}{0.38\textwidth}
      \begin{block}{Acc@1}
        \LARGE\bfseries 87.81 $\rightarrow$ 88.60\\
        \normalsize +0.79 điểm
      \end{block}
      \begin{block}{Acc@task}
        \LARGE\bfseries 88.01 $\rightarrow$ 88.32\\
        \normalsize +0.31 điểm
      \end{block}
      \begin{block}{Loss}
        \LARGE\bfseries 0.5228 $\rightarrow$ 0.4607\\
        \normalsize giảm khoảng 11.9\%
      \end{block}
    \end{column}
  \end{columns}
\end{frame}

% 35
\begin{frame}{Kết quả tốt hơn nhưng vẫn có trade-off}
  \begin{columns}[T,onlytextwidth]
    \begin{column}{0.48\textwidth}
      {\large\bfseries Điều đã cải thiện}\par\vspace{0.25cm}
      \begin{itemize}\setlength\itemsep{0.3cm}
        \item Final Acc@1 và Acc@task đều tăng.
        \item Loss giảm rõ, cho thấy classifier tự tin và phù hợp hơn.
        \item Energy weighting tốt hơn uniform ở cả accuracy và forgetting so với bước ngay trước.
      \end{itemize}
    \end{column}
    \begin{column}{0.48\textwidth}
      {\large\bfseries Điều chưa giải quyết}\par\vspace{0.25cm}
      \begin{itemize}\setlength\itemsep{0.3cm}
        \item Forgetting tăng từ 3.5889 lên 4.20: bản cuối quên task cũ nhiều hơn baseline.
        \item Backward giảm từ $-3.2556$ xuống $-3.9444$: tác động lên task cũ kém hơn.
        \item Một seed chưa chứng minh tính ổn định thống kê.
      \end{itemize}
    \end{column}
  \end{columns}
  \remember{Cấu hình hiện tại tối ưu final accuracy tốt hơn, nhưng chưa phải nghiệm tối ưu cho stability--plasticity trade-off.}
\end{frame}

% 36
\begin{frame}{Những thử nghiệm không được giữ lại và bài học}
  \small
  \begin{tabularx}{\textwidth}{>{\bfseries}p{3.25cm} X}
    \toprule
    Thử nghiệm & Kết luận \\
    \midrule
    Semantic top-k/gate/projection & Có ích nhẹ ở một vài metric nhưng không tăng Acc@1 ổn định. \\
    CFS distribution filter & Lọc quá chặt làm mất diversity, accuracy giảm. \\
    CFS mean initialization & Không tốt hơn khởi tạo ngẫu nhiên. \\
    CTIRD $K=3$ & Thiếu source task, kém $K=5$. \\
    Task temperature 0.2 & Kém temperature 0.1. \\
    CRCT 4 epochs & Có dấu hiệu over-correction, kém 3 epochs. \\
    Class-balanced replay batches & Acc@5 có thể tăng nhưng Acc@1 giảm. \\
    CTIRD con 0.15/0.25 & Đều kém con 0.20. \\
    \bottomrule
  \end{tabularx}
  \remember{Ablation âm cũng quan trọng: nó cho biết cải tiến nào thực sự đóng góp và tránh kể câu chuyện chỉ dựa trên lần chạy tốt.}
\end{frame}

% 37
\begin{frame}{Hạn chế và kế hoạch tiếp theo}
  \begin{enumerate}\setlength\itemsep{0.35cm}
    \item \textbf{Lặp lại baseline và cấu hình cuối}: xác nhận hai kết quả không phụ thuộc một lần chạy ngẫu nhiên.
    \item \textbf{Chạy nhiều seed}: ít nhất 3 seed để báo mean $\pm$ std và kiểm tra mức tăng 0.79 có ổn định không.
    \item \textbf{Benchmark thứ hai}: ImageNet-R hoặc 5-Datasets để kiểm tra khả năng tổng quát.
    \item \textbf{Giảm forgetting}: thử schedule/regularization CRCT mà không làm mất Acc@1 88.60.
    \item \textbf{Đo chi phí}: thời gian huấn luyện CFS, VRAM và số tham số MLP tạm thời.
    \item \textbf{Semantic thế hệ sau}: học alignment bằng visual training feature thay vì chỉ dựa class name.
  \end{enumerate}
  \remember{Kết quả hiện tại là bằng chứng thực nghiệm tốt cho một seed, chưa phải kết luận cuối cùng của nghiên cứu.}
\end{frame}

% 38
\begin{frame}{Tóm tắt toàn bộ công việc trong một trang}
  \small
  \begin{enumerate}\setlength\itemsep{0.22cm}
    \item HRM-PET gốc dùng CTIRD để giữ quan hệ liên task và Gaussian replay + CRCT để sửa classifier.
    \item Gaussian replay chưa chọn lọc; CRCT chưa dùng hết replay; code CTIRD legacy chọn source task chưa đúng mức liên quan.
    \item Ta chuyển ý tưởng CFS sang feature replay: học MLP contrastive theo lớp và chọn candidate ít giống nhau.
    \item Ta bật full replay và dành 25\% replay cho mẫu gần decision boundary.
    \item Ta thay class-level bottom-logit selection bằng top-5 task energy và weight KL theo energy nhưng giữ tổng lực CTIRD.
    \item Semantic-aware projection/gating đã được triển khai và ablation, nhưng không giữ vì chưa tăng Acc@1 ổn định.
    \item So với HRM-PET gốc cùng máy: Acc@1 $87.81\rightarrow88.60$, Acc@task $88.01\rightarrow88.32$, loss $0.5228\rightarrow0.4607$.
    \item Hạn chế còn lại: một seed; forgetting $3.5889\rightarrow4.20$ và Backward $-3.2556\rightarrow-3.9444$ kém hơn baseline.
  \end{enumerate}
  \remember{Thông điệp chính: cải thiện chất lượng/sử dụng replay và chọn đúng source task cho CTIRD tạo ra mức tăng đáng tin cậy nhất.}
\end{frame}

% 39
\appendix
\begin{frame}{Nguồn và file cần đọc lại}
  \small
  \begin{itemize}\setlength\itemsep{0.28cm}
    \item \textbf{Hybrid Re-matching for Continual Learning with Parameter-Efficient Tuning}: HRM-PET, DRM, CRM và CTIRD gốc.
    \item \textbf{Model Inversion with Layer-wise Optimization}: CFS và semantic-aware feature projection.
    \item \texttt{BAO\_CAO\_TOM\_TAT\_CAI\_TIEN\_HRM\_PET\_VI.md}: báo cáo tổng kết thay đổi và kết quả.
    \item \texttt{EXPERIMENT\_PROGRESS\_VI.md}: nhật ký từng ablation, tham số và kết luận.
    \item \texttt{CFS\_IMPLEMENTATION\_NOTE\_VI.md}: giải thích triển khai CFS.
    \item \texttt{SEMANTIC\_TOPK\_NOTE\_VI.md}: nhánh semantic và các biến thể.
    \item Code chính: \texttt{utils.py}, \texttt{hide\_tii\_engine.py}, \texttt{hrm\_lora\_wtp\_and\_tap\_engine.py}.
  \end{itemize}
\end{frame}

\end{document}
